\section{The surface and its orbifold multiplicities}\label{sec:definitions}

The construction in Theorem~\ref{thm:example} gives two projections:
$a:X\to E$ will be the quasi-Albanese map, while $j:Y\to\PP^1$ is
the logarithmic Iitaka fibration. We first describe their fibres.
We then set out the orbifold conventions needed to prove specialness
on every allowed upper model. All varieties and morphisms below are
algebraic over $\C$.

\subsection{The surface and its two projections}
We use the notation of Theorem~\ref{thm:example} and put
\[
 S=E\times\PP^1,\qquad H_i=E\times\{q_i\},\qquad
 A_i=b^{-1}(e_i,q_i).
\]
Write
\[
 \bar a=\operatorname{pr}_E\circ b:Y\to E,\qquad
 a=\bar a|_X,\qquad
 j=\operatorname{pr}_{\PP^1}\circ b:Y\to\PP^1.
\]
The points $e_i$ are permitted to coincide. The distinctness of the $q_i$
ensures that all blowup centres are distinct and that the curves $D_i$
are pairwise disjoint. We have
\begin{equation}\label{eq:intersection}
 A_i\simeq\PP^1,\qquad D_i\cdot A_j=\delta_{ij},\qquad
 A_i^\circ:=A_i\cap X\simeq\AA^1.
\end{equation}
Only the single intersection point with $D_i$ is removed from $A_i$.
These retained affine lines will account both for the orbifold
multiplicities and for the disappearance of puncture loops in $X$.

\paragraph{Fibres over the elliptic curve.}
For $e\in E$, put $I(e)=\{i:e_i=e\}$ and let $B_e$ be the strict
transform of $\{e\}\times\PP^1$. The scheme fibre of $\bar a$ is
\[
 B_e+\sum_{i\in I(e)}A_i,
\]
with all multiplicities one. Its components in the open surface are
\[
 B_e\cap X\simeq\PP^1\setminus\{q_j:j\notin I(e)\},
 \qquad A_i^\circ\simeq\AA^1\quad(i\in I(e)).
\]
Each $A_i^\circ$ meets $B_e\cap X$ once transversely, and the exceptional
components are mutually disjoint. The point of attachment represents the
vertical tangent direction at the original centre; the omitted point
represents the horizontal tangent direction and is different. Thus every
fibre of $a$ is connected. For general $e$, the set $I(e)$ is empty and
\begin{equation}\label{eq:general-a-fibre}
 a^{-1}(e)\simeq\PP^1\setminus\{q_1,\ldots,q_n\}.
\end{equation}

\paragraph{Elliptic fibres and the logarithmic canonical divisor.}
For $t\notin\{q_i\}$, the fibre of $j$ is an unchanged complete elliptic
curve contained in $X$. At $q_i$ its scheme fibre is
\begin{equation}\label{eq:j-fibre}
 j^*(q_i)=D_i+A_i,
\end{equation}
both components having multiplicity one. In particular, all fibres of
$j$ are connected and $j_*\OO_Y=\OO_{\PP^1}$.

The canonical divisor and the strict-transform formula give
\begin{align}
 K_Y&=b^*K_S+\sum_iA_i, &
 D&=b^*\Bigl(\sum_iH_i\Bigr)-\sum_iA_i,\nonumber\\
 K_Y+D&\sim j^*\OO_{\PP^1}(n-2).\label{eq:log-canonical}
\end{align}
Projection formula therefore gives
\[
 H^0\bigl(Y,m(K_Y+D)\bigr)
   \simeq H^0\bigl(\PP^1,\OO_{\PP^1}(m(n-2))\bigr)
 \qquad(m\geq1).
\]
Thus $\kappa(Y,K_Y+D)=1$, and $j$ is the Iitaka fibration: the
logarithmic pluricanonical sections determine $j$ up to an embedding of
its base. The complete elliptic fibres of $j$ will constrain every
possible curve fibration when we prove specialness.

\subsection{The infimum convention}
A smooth orbifold pair is a smooth variety $Y$
together with a simple normal crossing $\Q$-divisor
\[
 D=\sum_P\left(1-\frac1{m_P(D)}\right)P,
 \qquad m_P(D)\in\Z_{\geq1}\cup\{\infty\}.
\]
Only finitely many coefficients are nonzero. We use $1/\infty=0$;
thus a component of a reduced boundary has multiplicity $\infty$, whereas
a prime divisor outside its support has multiplicity $1$.

Here a \emph{fibration} is a surjective morphism with connected fibres.
Let $f:Y\to Z$ be such a morphism, with $Y$ smooth projective and $Z$
smooth. For a prime divisor $Q\subset Z$ and a prime divisor $P\subset Y$
dominating $Q$, write $r(f,P)$ for the coefficient of $P$ in $f^*Q$.
The orbifold base divisor is
\begin{equation}\label{eq:infimum}
 \Delta(f,D)=\sum_Q\left(1-\frac1{m_f(Q)}\right)Q,
 \qquad
 m_f(Q)=\inf_{f(P)=Q}r(f,P)m_P(D).
\end{equation}
This is the infimum convention of
\cite[Definition 2.1]{CDHP2026}; it is not a greatest-common-divisor
convention. A divisor $P$ with $\codim_Zf(P)\geq2$ is
\emph{$f$-exceptional} and does not occur in the infimum for the immediate
base. It may occur for a later composition.

For the surface above, \eqref{eq:j-fibre} gives
\begin{equation}\label{eq:j-base-zero}
 m_j(q_i)=\min\{1\cdot\infty,1\cdot1\}=1,
 \qquad \Delta(j,D)=0.
\end{equation}
The unchanged fibres give multiplicity one at every other point. Thus
$(Y,D)$ has logarithmic Kodaira dimension one, but the orbifold base of
$j$ is $(\PP^1,0)$. The retained component $A_i$ is what makes
its multiplicity at $q_i$ equal to one.

\subsection{Upper models and specialness}
A dominant morphism of smooth orbifold pairs $u:(Y',D')\to(Y,D)$ is an
\emph{orbifold morphism} if, for every prime divisor $P\subset Y$ and
every component $P'$ of $u^*P$, its coefficient $r$ in that pullback
satisfies
\[
 r\,m_{P'}(D')\geq m_P(D).
\]
It is an \emph{elementary birational orbifold morphism} if, in addition,
it is birational and $u_*D'=D$. In particular, when $D$ is reduced, every
component of $u^{-1}D$ must lie in the coefficient-one part of $D'$.
An equality of pushed-forward boundaries alone is weaker.

A fibration $f:(Y,D)\to Z$ is \emph{neat} if there is a commutative
birational diagram
\[
\begin{tikzcd}[column sep=large]
 (Y,D) \arrow[r,"f"] \arrow[d,"w"'] & Z \arrow[d,"v"]\\
 (Y_0,D_0) \arrow[r,"f_0"'] & Z_0
\end{tikzcd}
\]
with smooth spaces, $w$ an elementary birational orbifold morphism, and every $f$-exceptional divisor also $w$-exceptional. The
definition is the one in \cite[Definition 2.3]{CDHP2026}. A fibration is
\emph{high} if its morphism to the base equipped with the induced divisor
$\Delta(f,D)$ is an orbifold morphism. Birational equivalence of
fibrations is generated by diagrams in which the source modification is
an elementary birational orbifold morphism and the base modification is
birational. The \emph{canonical dimension} of a fibration is the infimum
of the Kodaira dimensions of its orbifold bases over these models. Here
the Kodaira dimension of a $\Q$-divisor is the exponent measuring growth
of its plurisections, with value $-\infty$ if all positive plurisection
spaces vanish. For a neat fibration,
\begin{equation}\label{eq:neat-kappa}
 \kappa(f,D)=\kappa\bigl(Z,K_Z+\Delta(f,D)\bigr)
\end{equation}
by \cite[Lemma 2.7]{CDHP2026}, attributed there to
\cite[Corollaire 5.11(3)]{Campana2011}.

\begin{definition}\label{def:special}
A smooth projective orbifold pair $(Y,D)$ is special if for every
elementary birational orbifold morphism
$u:(Y',D')\to(Y,D)$ with $(Y',D')$ smooth projective and every neat fibration
$g:(Y',D')\to Z$ with $\dim Z>0$,
\[
 \kappa\bigl(Z,K_Z+\Delta(g,D')\bigr)<\dim Z.
\]
A smooth open variety $X=Y\setminus D$ with reduced boundary is
logarithmically special when its smooth projective compactification pair
is special in this sense.
\end{definition}

For a smooth projective curve $C$, a divisor $K_C+\Delta$ is of general
type precisely when its degree is positive. In dimension two we will
also use the elementary fact that a birational fibration is equivalent
to the identity fibration on its source. Consequently its canonical
dimension is no larger than the logarithmic Kodaira dimension of that
source. Equation~\eqref{eq:neat-kappa} then makes this useful for
Definition~\ref{def:special}.

\paragraph{Comparison with the open-variety convention.}
The condition above implies specialness in the sense of
\cite[Definition 2.1]{CDY2025}, which is the convention used by the
external inputs in Section~\ref{sec:linear}. To see the needed
implication, start with an open fibration and proper birational
modification allowed there. Compactify the modification over $(Y,D)$,
resolve the graph of the fibration, and choose a neat upper model as
in \cite[Proposition 2.8]{CDHP2026}. On this model, let $D_{\mathrm{open}}$
be the reduced inverse image of the original boundary and let
$D_{\mathrm{high}}$ be the boundary of the upper orbifold pair.
The orbifold-morphism condition over the reduced divisor $D$ forces
$D_{\mathrm{open}}\leq D_{\mathrm{high}}$. On the same resolved
morphism $g$, formula~\eqref{eq:infimum} gives
\[
 \Delta(g,D_{\mathrm{open}})\leq\Delta(g,D_{\mathrm{high}}).
\]
Consequently the Kodaira dimension of the former orbifold base is no
larger than that of the latter, which is smaller than the dimension
of the base by Definition~\ref{def:special}. The infimum over models
in the open-variety definition is at most this value. This proves
the implication used here; no identification of the two classes of
models is needed.
